\section{Boundary tests from rare pooled collisions}
\label{sec:rare-stationary}
The preceding stationary construction estimates a difference of two
collision probabilities to the accuracy of a small boundary mass.
Near an exposed boundary one can instead arrange that an exterior
test has a zero collision probability. An interior test then needs only
one observed collision. The compass direction remains unobserved.
To cover the directions at which two compass vertices maximize the same
support functional, we test a finite set of candidate vertices and take
the conjunction of their outcomes.

The compass length in this section is chosen once from the fixed
priors. Let $D_*$ be a known upper bound for the diameter of every
available launch footprint, and require
\begin{equation}\label{eq:rare-margin}
                         2t+D_*<d_0.
\end{equation}
For the calibrated one-scale experiment one may take $D_*=\diam K$;
for the two known scales of Section~\ref{sec:unknown-footprint}, take
$D_*=2\lambda_2R_K$. The extra factor of two accounts for both the
shift of a nominal center and the attempted flight. This sufficient
design condition is stronger than the one used for the earlier
reciprocal estimates. Those estimates remain valid under their stated
conditions. Whenever the footprint diameter bound is strictly below
$d_0$, a positive dyadic $t$ satisfying \eqref{eq:rare-margin} can be
chosen independently of the requested accuracy. No new output or
direction control is required.

\subsection{Coarse normal information}
The expanded components $P=C+(-K)$ have a uniform positive lower
curvature radius, a uniform support regularity bound, and a positive
separation. Fix a collar width $d_c>0$ smaller than their lower rolling
radius and such that
\begin{equation}\label{eq:rare-collar}
            d_c<t/8,\qquad D_*+2t+d_c<d_0.
\end{equation}
Reducing $d_c$ below a fixed fraction of these bounds will be harmless.
For every $y$ in this collar of a complete expanded component, there
are a boundary point $y_0$, its outward unit normal $n$, and a signed
depth $d\in[-d_c,d_c]$ such that
\begin{equation}\label{eq:rare-signed-depth}
              y=y_0-dn,\qquad n\cdot y=p_P^{\rm lab}(n)-d.
\end{equation}
Here $d>0$ inside $P$. The representation is unique in the chosen
collar: the support function $p_P-d$ still has a positive curvature
radius for every inward offset in this range, and describes the inner
parallel boundary. Exterior metric projections are unique by convexity.

\begin{lemma}[Fixed-accuracy normals for radial brackets]
\label{lem:rare-coarse-normals}
At any prescribed confidence $1-\varepsilon_c$, a prior-dependent
finite number of the earlier reciprocal queries supplies an interior
center for each protected expanded component and isolated radial
brackets containing its boundary in every radial direction. Each bracket
has a tubular neighborhood of one fixed positive prior-controlled radius
lying in the collar \eqref{eq:rare-signed-depth}. For each radial
direction the same finite data compute a rational vector $\widehat n$
satisfying
\begin{equation}\label{eq:rare-normal-error}
                         |\widehat n-n|\le1/32
\end{equation}
for the projection normal of every point in that tubular neighborhood.
The number
of attempts is at most $C\log(C/\varepsilon_c)$. For unknown
footprints the construction uses only their uniform priors.
\end{lemma}
\begin{proof}
Apply the fixed-layer version of Proposition~\ref{prop:jitter-query}
and the coarse acquisition of Lemma~\ref{lem:coarse} to the expanded
components, obtaining interior centers and isolated radial intervals.
With respect to each center, take a fixed angular mesh of width $h_c$
and refine its radial values to error $e_c\le c h_c^s$ by the same
reciprocal queries. All these parameters are fixed from the priors.
Lemma~\ref{lem:radial} yields a finite smooth radial function
$\widetilde R$ with
\[
 \|\widetilde R-R_P\|_\infty\le C h_c^s=:E_c,
 \qquad
 \|\widetilde R-R_P\|_{C^1}\le C h_c^{s-1}.
\]
Choose $h_c$ small enough that both error bounds meet the fixed
reserves below. In radial direction $u(\varphi)$, use a certified
interval containing
$[\widetilde R(\varphi)-2E_c,\widetilde R(\varphi)+2E_c]$.
Its endpoints enclose the true radial boundary, and its width can be
made arbitrarily small at this fixed stage. It remains in the original
isolated interval after reducing $h_c$ and its numerical reserve.

The normal to a positive radial graph is
\[
 \frac{R_P(\varphi)u(\varphi)
                 -R_P'(\varphi)u^\perp(\varphi)}
        {\sqrt{R_P(\varphi)^2+R_P'(\varphi)^2}}.
\]
The inball and regularity bounds make this expression uniformly
Lipschitz in $R_P,R_P'$. Substitution of $\widetilde R$ therefore
approximates the true radial boundary normal to any prescribed fixed
accuracy. The projection normals throughout the bracket have the same
property. To verify this last assertion quantitatively, let $r_*>0$
be a uniform lower rolling radius of $P$. Write $P=Q+r_*\overline B$,
using the support identity from Section~\ref{sec:stationary}. If
$x_i\in\partial P$ has outer normal $n_i$, the monotonicity of the
support-point relation for $Q$ gives
\[
 (x_1-x_2)\cdot(n_1-n_2)\ge r_*|n_1-n_2|^2,
 \qquad |n_1-n_2|\le |x_1-x_2|/r_*.
\]
The projection of any collar point lies within twice its distance
from the true radial boundary point. Hence on a tubular neighborhood
of radius $a_c$ about the bracket its normal differs from the radial
boundary normal by $O(E_c+a_c)$. Choose $h_c$ and then a fixed
$a_c>0$ so small that the whole neighborhood lies in the chosen collar
and that this error and the radial approximation error sum to less
than $1/64$. Compute a rational normal approximation with another
error less than $1/64$. This proves \eqref{eq:rare-normal-error}
uniformly on that neighborhood; its radius is independent of all
subsequent accuracy parameters.

Only a fixed number of radial values, bisections and reciprocal command
centers was used. All requested probability accuracies are fixed, so
their conditional confidence allocation costs
$C\log(C/\varepsilon_c)$ attempts. Every step precedes footprint
subtraction and depends only on the uniform bounds for the expanded
bodies. It consequently applies at either unknown-footprint scale.
\end{proof}

\subsection{A test using only positive collision records}
Put $E=\{e_1,-e_1,e_2,-e_2\}$ and, for the rational vector supplied
by Lemma~\ref{lem:rare-coarse-normals}, form the finite set
\begin{equation}\label{eq:rare-candidates}
 V(\widehat n)=
 \left\{v\in E:\ \widehat n\cdot v
       \ge\max_{w\in E}\widehat n\cdot w-\tfrac14\right\}.
\end{equation}
All these comparisons are rational. This set contains every maximizer
of $n\cdot v$ over $E$: the two normal errors cost at most $1/16$.
Moreover, every candidate satisfies
\begin{equation}\label{eq:rare-positive-normal}
 n\cdot v\ge\frac1{\sqrt2}-\frac14-\frac1{16}>\frac38.
\end{equation}
The threshold in \eqref{eq:rare-candidates} is positive, so the set
contains at most two compass vertices. In particular the construction
does not require a unique maximizing direction.

\begin{proposition}[A rare-collision boundary query]
\label{prop:rare-query}
Assume \eqref{eq:rare-margin} and condition on the successful coarse
event of Lemma~\ref{lem:rare-coarse-normals}. For a target $y$ in one
of its bracket neighborhoods, a scale $0<e<d_c$ and a conditional failure
allowance $0<\varepsilon<1/4$, pooled forward commands at at most two
shifted centers return the correct membership label whenever $y$ is
outside the expanded component or is inside it at distance at least
$e$ from its boundary. The label in the intervening inner layer is
unrestricted. The query uses at most
\begin{equation}\label{eq:rare-query-cost}
                     C e^{-(\gamma+3/2)}\log(C/\varepsilon)
\end{equation}
attempted collision bits. An exterior target is labelled zero with
probability one. No launch position, free-start indicator or compass
direction is observed.
\end{proposition}
\begin{proof}
For every $v\in V(\widehat n)$ make fresh pooled forward preparations
at $x_v=y+tv$. Mark this candidate positive if its batch contains at
least one collision. Return one precisely when every candidate is
marked positive. The random direction in each preparation still has
the prescribed uniform law on $tE$.

First exclude all other obstacles from these preparations. At the
projection in \eqref{eq:rare-signed-depth}, support addition writes
$y_0=c_0-z_0$, where $c_0\in C$ and $z_0\in K$ are the support
points with normals $n$ and $-n$. Every point of every attempted
segment has the form
\[
 c_0+(z-z_0)+tv+utw-dn,
       \qquad z\in K,\quad w\in E,\quad 0\le u\le1.
\]
Its distance from $c_0$ is at most $D_*+2t+d_c<d_0$. The physical
separation therefore excludes every component other than $C$, for
all launch offsets and all compass directions.

Suppose that $y$ is exterior, so $d<0$, and choose a true maximizing
vertex $v_*\in V(\widehat n)$. For every $z\in K$, $w\in E$ and
$0\le u\le1$, support addition gives
\begin{align*}
 n\cdot(y+tv_*+z+utw)
 &\ge p_C^{\rm lab}(n)-d+t(n\cdot v_*+u n\cdot w)\\
 &>p_C^{\rm lab}(n).
\end{align*}
The second line uses symmetry of $E$ and maximality of $v_*$, which
imply $n\cdot v_*+u n\cdot w\ge0$. Thus every segment for this
candidate misses $C$, and the preceding separation excludes all other
collisions. Its collision probability is zero. Its mark, and therefore
the conjunction, is deterministically zero. Possible boundary tangencies
are irrelevant to the asserted exterior and positive-depth labels.

Now suppose that $y$ has inner depth $d\ge e$. By
Lemma~\ref{lem:cap-mass},
\[
        \Prob\{y+Z\in C\}=v_j(y)
                    \ge c e^{\gamma+3/2}=:b_e.
\]
Decrease the fixed constant $c$ if necessary so that $b_e\le1$.
For every candidate $v$ and every $z\in K$,
\[
 n\cdot(y+tv+z)
       \ge p_C^{\rm lab}(n)-d+t n\cdot v
       >p_C^{\rm lab}(n),
\]
by \eqref{eq:rare-collar} and \eqref{eq:rare-positive-normal}.
All starts are consequently free. On the independent events
$y+Z\in C$ and $a=-tv$, the attempted segment ends in $C$ and
collides. The pooled success probability for each candidate is at
least $b_e/4$; the selected direction need never be revealed.

Use
\[
            M_e=\left\lceil\frac4{b_e}
                                     \log\frac2\varepsilon\right\rceil
\]
fresh preparations for each candidate. Its probability of having no
collision is at most $\exp(-M_eb_e/4)\le\varepsilon/2$.
A union bound over at most two candidates proves correctness at
positive depth and \eqref{eq:rare-query-cost}. Fresh preparations give
the same bound conditional on every preceding adaptive history.
No upper bound or regularity assumption on the density was used.
\end{proof}

\subsection{The improved stationary reconstruction rate}
\begin{theorem}[Stationary reconstruction by rare collisions]
\label{thm:rare-stationary}
Fix the priors of Theorem~\ref{thm:stationary-jitter} and choose a
fixed compass length satisfying $2t+\diam K<d_0$. There is a finite
adaptive controller, independent of the stationary density, whose
geometric loss is at most $C\nu$ and whose primitive discrete data
are correct with probability at least $1-\delta$, uniformly over the
same physical and density classes. For sufficiently small $\nu>0$
and $0<\delta<1/4$, its costs satisfy
\begin{align}
 J_\nu&\le C\nu^{-1/(s-2)}\log(C/\nu),
                                      \label{eq:rare-centers}\\
 N_\nu&\le C\nu^{-((\gamma+3/2)s+1)/(s-2)}
              \log(C/\nu)\log(C/(\nu\delta)).
                                      \label{eq:rare-attempts}
\end{align}
Every preparation is counted. The launch footprint and compass length
stay fixed, and all nominal centers and physical launches lie in a
prior-dependent fixed aperture. Dyadic nominal positioning to
$O(\nu^{s/(s-2)})$ suffices.

For the two known scales and unknown footprint of
Theorem~\ref{thm:self-finite}, the same conclusions and costs hold
under $2t+2\lambda_2R_K<d_0$, including laboratory $C^2$ footprint
error at most $C\nu$. No footprint-function or density evaluations
are then required. The exact scales, common dilation origin and fixed
geometric and boundary-mass priors retain their stated roles.
\end{theorem}
\begin{proof}
Perform the fixed-accuracy construction of
Lemma~\ref{lem:rare-coarse-normals}, reserving failure probability
$\delta/2$. Put $h=2\pi/m\asymp\nu^{1/(s-2)}$ and $e=c h^s$.
At each of the $m$ equally spaced radial angles, use
Proposition~\ref{prop:rare-query} for the midpoint tests of the
corresponding coarse bracket. Its rational candidate vector is fixed
for that radial direction and remains valid throughout the bracket.
The relaxed bisection invariant of Lemma~\ref{lem:bisection} applies
to the inner indeterminate layer as well. The inball converts a normal
depth of order $e$ to a radial uncertainty of order $e$. Thus
$O(\log(C/e))$ tests per direction give radial values with error
$O(e)$, without a monotonicity assumption on erroneous layer labels.

Lemma~\ref{lem:radial} gives expanded bodies $\widehat P_C$ with
support error $O(h^s)$ in supremum norm and $O(h^{s-2})$ in $C^2$.
For a known footprint, subtract $p_{-K}^{\rm lab}$ and apply the
curvature, final smoothing and period-locking arguments of
Theorem~\ref{thm:stationary-jitter}. These deterministic steps give
the stated physical output and retain the primitive conclusions.

There are at most $Q_h=C(1+h^{-1}\log(C/h))$ fine queries. Give
each one conditional failure allowance $\delta/(2Q_h)$. On the
successful coarse event, Proposition~\ref{prop:rare-query} and a
union bound give simultaneous valid labels with probability at least
$1-\delta/2$. The attempts are bounded by
\[
 C\log(C/\delta)
   +C h^{-1}e^{-(\gamma+3/2)}
                   \log(C/h)\log(C/(h\delta)).
\]
Substituting $e=c h^s$ proves \eqref{eq:rare-attempts}. At most two
shifted command centers replace each fine target; coarse reciprocal
stencils have fixed size. This proves \eqref{eq:rare-centers}.
An enlargement of the earlier fixed aperture by the fixed shifts
protects all these commands and their launch supports.

For numerical commands, round each target once to a dyadic grid of
mesh $O(e)$ and then add the exact dyadic vector $tv$. For sufficiently
small $\nu$, every rounded target lies in the fixed tubular neighborhood
of its radial bracket. The same candidate vector therefore satisfies
\eqref{eq:rare-normal-error}, even though rounding may move the target
off its radial line. Rounding changes distances to the expanded component
by $O(e)$ and hence only enlarges the indeterminate layer by that
amount. It does not turn a claimed exterior test at the rounded point
into a probability estimate. Conservative dyadic lower bounds for
$b_e$ and integer batch counts give the same costs. Their descriptions,
and those of the centers, have
$O(\log(C/\nu)+\log(C/\delta))$ digits per batch.

In the two-scale experiment apply this construction separately to
$P_{i,C}=C+(-\lambda_iK)$, using
\eqref{eq:scaled-density-bound} and the common diameter envelope
$2\lambda_2R_K$. The stronger separation assumption gives
\eqref{eq:rare-margin} at both scales. Match complete components by
Lemma~\ref{lem:scale-matching}, and use the averaged linear support
inverse and convexity verification of Theorem~\ref{thm:self-finite}.
Both scales have the same exponents and uniform prior-controlled
constants. Splitting the confidence budget and adding their costs
therefore preserves \eqref{eq:rare-centers}--\eqref{eq:rare-attempts}.
\end{proof}

For $\gamma=0$, the new upper power is $(3s/2+1)/(s-2)$.
The lower bound of Theorem~\ref{thm:stationary-lower}, specialized to
a fixed sufficiently small uniform-disk footprint and the present
compass choice, consequently gives, on a fixed physical class containing
its packing,
\begin{equation}\label{eq:rare-stationary-bracket}
 c\nu^{-(s+1)/(s-2)}
       \le N^*_{\rm stat}(\nu)
       \le C\nu^{-(3s/2+1)/(s-2)}\log^2(C/\nu).
\end{equation}
The difference between the displayed exponents is now
$s/(2(s-2))$. These bounds do not assert an exact minimax rate.
The improvement comes from a rare positive event whose absence has
an exact geometric explanation on the exterior side; it does not
follow from increasing the precision of reciprocal mean cancellation.
